\chapter[Transformacions naturals]{\texorpdfstring{Transformacions naturals\\ i equivalències}{Transformacions naturals i equivalències}}
\label{cap:transformacions}

\begin{objectius}
\apartat{Prerequisits} Functors, els homfunctors i el bifunctor $\Hom$ (cap.~\ref{cap:functors}). Per a alguns exemples convé, a més, una familiaritat bàsica amb els espais vectorials i, en particular, amb l'espai dual $V^*$ i el doble dual $V^{**}$.
\apartat{Objectius} En acabar el capítol, el lector ha de saber: definir una transformació natural i verificar la naturalitat via quadrats commutatius; reconèixer isomorfismes naturals; treballar amb la categoria de functors $\cat{D}^{\cat{C}}$; compondre transformacions horitzontalment; i caracteritzar una equivalència de categories com un functor ple, fidel i essencialment exhaustiu.
\end{objectius}

\section{Cap al concepte de transformació natural}\label{sec:cap-transformacio}

Al capítol anterior vam veure que els functors són els morfismes entre categories: relacionen dues categories respectant-ne tota l'estructura. Ara bé, sovint tenim \emph{dos} functors $F,G\colon\cat{C}\to\cat{D}$ que van de la mateixa categoria a la mateixa categoria, i volem comparar-los. La pregunta natural és, doncs, si hi ha una noció de \emph{morfisme entre functors}.

Aquesta idea ---comparar functors de manera coherent--- va ser històricament una de les motivacions originals per introduir les categories. De fet, es diu sovint que les categories es van definir per poder definir els functors, i els functors per poder definir les transformacions naturals. El mot \enquote{natural} té aquí un sentit tècnic precís que farem explícit: una comparació és natural quan fa commutar els quadrats corresponents a tots els morfismes. Sovint s'expressa dient que \enquote{no depèn de cap elecció arbitrària}; convé retenir, però, que es tracta d'una intuïció orientadora i no d'una definició: una construcció pot haver requerit eleccions i, tanmateix, resultar natural. La condició precisa és sempre la commutativitat dels quadrats.

Abans de formalitzar-ho, val la pena veure un exemple que fa palès aquest \enquote{principi de coherència}, aquest cop entre \emph{sintaxi} i \emph{semàntica}. Treballem amb la lògica proposicional. Per a cada conjunt $X$ de variables proposicionals, considerem dues construccions:
\[
F(X)=\{\text{fórmules proposicionals amb variables de }X\}
\]
i
\[
G(X)=\{\text{funcions de veritat sobre }X\},
\]
on una \emph{funció de veritat} sobre $X$ és una aplicació que assigna un valor de $\{0,1\}$ a cada valoració $v\colon X\to\{0,1\}$ de les variables. Totes dues són functorials en $X$: un canvi de variables transforma les fórmules per substitució i les funcions de veritat per transport de les valoracions. A cada conjunt de variables $X$ li associem, a més, l'aplicació que interpreta cada fórmula com la seva taula de veritat:
\[
\eta_X\colon F(X)\longrightarrow G(X),
\qquad
\eta_X(\varphi)(v)=
\begin{cases}
1 & \text{si }v\models\varphi,\\
0 & \text{si }v\not\models\varphi.
\end{cases}
\]
El fet clau és que aquesta interpretació és \emph{compatible amb el canvi de variables}: donat un canvi de variables $f\colon X\to Y$, substituir primer i interpretar després dona la mateixa funció de veritat que interpretar primer i transportar després. Per exemple, de $\varphi=p\land q$ i la substitució $p\mapsto r,\ q\mapsto\neg r$ obtenim $r\land\neg r$, i tant se val si calculem la taula de veritat abans o després de substituir: el resultat és la mateixa funció (aquí, la funció constant $0$). Aquesta compatibilitat, per a tot canvi de variables alhora, és el que a la secció següent anomenarem la \textbf{naturalitat} de la família $(\eta_X)$:
\[
\boxed{\text{substituir i després interpretar} \;=\; \text{interpretar i després canviar les variables}.}
\]

Val la pena notar una diferència amb els exemples \enquote{d'identificació canònica} que veurem després (com el doble dual d'un espai vectorial): aquí $\eta_X$ \emph{no} és injectiva ---dues fórmules diferents poden tenir la mateixa taula de veritat, com $p\lor q$ i $q\lor p$---, de manera que no estableix cap isomorfisme entre sintaxi i semàntica. El que expressa la naturalitat no és una identificació, sinó que \emph{el significat d'una fórmula no depèn de com hàgim anomenat les variables}. La mateixa idea s'estén, amb més maquinària, a la lògica de primer ordre, on les fórmules s'interpreten com els predicats que defineixen sobre una estructura fixada; ho veurem desplegat com a exemple formal a l'exemple~\ref{ex:interpretacio-primer-ordre} de la secció~\ref{sec:exemples-transf}.

\section{Definició de transformació natural}

Siguin $F,G\colon\cat{C}\to\cat{D}$ dos functors. Volem, per a cada objecte $A$ de $\cat{C}$, un morfisme de $\cat{D}$ que vagi de $F(A)$ a $G(A)$; i volem que aquests morfismes siguin compatibles amb l'acció de $F$ i $G$ sobre les fletxes.

\begin{definicio}\index{transformació natural}
\label{def:transf-natural}
Una \textbf{transformació natural} $\eta\colon F\Rightarrow G$ entre dos functors $F,G\colon\cat{C}\to\cat{D}$ assigna, a cada objecte $A$ de $\cat{C}$, un morfisme de $\cat{D}$
\[
\eta_A\colon F(A)\to G(A),
\]
anomenat \textbf{component} de $\eta$ a $A$, de manera que, per a cada morfisme $f\colon A\to B$ de $\cat{C}$, el quadrat següent commuta:
\[
\begin{tikzpicture}[baseline=(m.center)]
  \matrix (m) [matrix of math nodes, row sep=2.8em, column sep=3.4em] {
    F(A) & G(A) \\
    F(B) & G(B) \\
  };
  \draw[-{Stealth[scale=1.1]}] (m-1-1) -- node[above] {$\eta_A$} (m-1-2);
  \draw[-{Stealth[scale=1.1]}] (m-1-1) -- node[left] {$F(f)$} (m-2-1);
  \draw[-{Stealth[scale=1.1]}] (m-1-2) -- node[right] {$G(f)$} (m-2-2);
  \draw[-{Stealth[scale=1.1]}] (m-2-1) -- node[below] {$\eta_B$} (m-2-2);
\end{tikzpicture}
\]
és a dir,
\[
G(f)\comp\eta_A=\eta_B\comp F(f).
\]
Aquesta igualtat s'anomena \textbf{condició de naturalitat}, i el quadrat, \textbf{quadrat de naturalitat}.
\end{definicio}

\begin{observacio}
La condició de naturalitat és el cor de la definició. Diu que \enquote{transformar i després aplicar $G(f)$} coincideix amb \enquote{aplicar $F(f)$ i després transformar}. Els components $\eta_A$ es trien un per a cada objecte, però no lliurement: han d'encaixar tots alhora amb l'acció dels functors sobre els morfismes. Convé precisar bé què exigeix aquesta condició i què no: la naturalitat demana que els components escollits siguin compatibles amb tots els morfismes; no demana que siguin l'única elecció possible. És una condició de compatibilitat, no d'unicitat ni d'absència de paràmetres.

Un exemple ho fa evident. Sobre els espais vectorials reals, fixem un escalar $\lambda\in\mathbb{R}$ i prenem, per a cada espai $V$, la component
\[
\eta_V=\lambda\,\id_V\colon V\to V.
\]
Aquesta família és una transformació natural de la identitat en ella mateixa, perquè per a tota aplicació lineal $T\colon V\to W$,
\[
T\comp(\lambda\,\id_V)=\lambda\,T=(\lambda\,\id_W)\comp T,
\]
de manera que el quadrat de naturalitat commuta; quan $\lambda\neq 0$, és fins i tot un isomorfisme natural. Tenim, doncs, tota una família de transformacions naturals, una per cada valor de $\lambda$: la naturalitat es compleix per a qualsevol $\lambda$ fixat, tot i que la construcció depèn d'aquesta elecció. Naturalitat no vol dir, per tant, \enquote{sense eleccions}, sinó \enquote{compatible amb tots els morfismes}.
\end{observacio}

\begin{observacio}[Una família que \emph{no} és natural]\index{transformació natural!contraexemple}
És igual d'instructiu veure una família de components que \emph{falla} la condició de naturalitat. Treballem a $\cat{Set}$ amb $F=G=\id_{\cat{Set}}$, de manera que una transformació $\id\Rightarrow\id$ hauria de donar, per a cada conjunt $A$, una aplicació $\eta_A\colon A\to A$ que commuti amb \emph{totes} les aplicacions. Provem a triar, per a cada conjunt no buit $A$, un element distingit $a_0\in A$ i posem $\eta_A$ igual a l'aplicació constant de valor $a_0$. El quadrat de naturalitat per a una aplicació $f\colon A\to B$ demanaria $f\comp\eta_A=\eta_B\comp f$, és a dir $f(a_0)=b_0$ per a tota $f$.

Aquesta condició falla. Prenem $A=B=\{0,1\}$ amb la mateixa tria $a_0=b_0=0$, i l'aplicació constant $f\equiv 1$. El quadrat
\[
\begin{tikzpicture}[baseline=(m.center)]
  \matrix (m) [matrix of math nodes, row sep=2.8em, column sep=3.8em] {
    A & A \\
    B & B \\
  };
  \draw[-{Stealth[scale=1.1]}] (m-1-1) -- node[above] {$\eta_A$} (m-1-2);
  \draw[-{Stealth[scale=1.1]}] (m-1-1) -- node[left] {$f$} (m-2-1);
  \draw[-{Stealth[scale=1.1]}] (m-1-2) -- node[right] {$f$} (m-2-2);
  \draw[-{Stealth[scale=1.1]}] (m-2-1) -- node[below] {$\eta_B$} (m-2-2);
\end{tikzpicture}
\]
no commuta: seguint el camí de dalt i a la dreta, $f(\eta_A(x))=f(0)=1$; seguint el de l'esquerra i a baix, $\eta_B(f(x))=\eta_B(1)=0$. Com que $1\neq 0$, els dos camins difereixen, i la família $(\eta_A)$ \emph{no} és una transformació natural. El problema no és la tria en si ---a l'exemple anterior també triàvem un $\lambda$---, sinó que aquesta tria \emph{no és compatible} amb totes les fletxes: escollir un element distingit a cada conjunt no respecta les aplicacions que no el preserven. La naturalitat és exactament aquesta compatibilitat global, i aquí es trenca.
\end{observacio}

\section{Exemples de transformacions naturals}\label{sec:exemples-transf}

\begin{exemple}[Transformació identitat]\index{transformació natural!identitat}
Per a tot functor $F\colon\cat{C}\to\cat{D}$ hi ha la \textbf{transformació natural identitat} $\id_F\colon F\Rightarrow F$, amb components $(\id_F)_A=\id_{F(A)}$. El quadrat de naturalitat commuta trivialment, perquè els dos camins són $F(f)$.
\end{exemple}

\begin{exemple}[Preordres]\index{transformació natural!entre preordres}
Siguin $P,Q$ dos preordres i $F,G\colon P\to Q$ dos functors, és a dir, dues aplicacions monòtones. Una transformació natural $\eta\colon F\Rightarrow G$ dona, per a cada $x\in P$, un morfisme $\eta_x\colon F(x)\to G(x)$; però a $Q$ hi ha una fletxa $F(x)\to G(x)$ exactament quan $F(x)\leq G(x)$. La condició de naturalitat es compleix automàticament, perquè entre dos objectes d'un preordre hi ha com a màxim una fletxa. Per tant,
\[
\text{existeix }\eta\colon F\Rightarrow G
\quad\Longleftrightarrow\quad
F(x)\leq G(x)\ \text{per a tot }x.
\]
En el llenguatge de l'ordre, una transformació natural entre aplicacions monòtones és exactament la relació \enquote{$F$ és puntualment menor o igual que $G$}.
\end{exemple}

\begin{exemple}[Interpretació de fórmules de primer ordre]\label{ex:interpretacio-primer-ordre}\index{transformació natural!interpretació lògica}\index{lema de substitució}
Reprenem el fil de la lògica que hem obert a la secció~\ref{sec:cap-transformacio} (\enquote{Cap al concepte de transformació natural}), ara en primer ordre i amb la naturalitat plantejada com a exemple formal. Fixem un llenguatge de primer ordre i una estructura $M$. Per a cada conjunt finit de variables $X$ ---un \emph{context}--- considerem dues construccions:
\[
F(X)=\{\text{fórmules amb variables lliures entre }X\}
\]
i
\[
G(X)=\{\text{predicats definibles sobre }X\},
\]
on un \emph{predicat definible} és un subconjunt de $M^X$ (el conjunt d'assignacions $s\colon X\to M$) de la forma
\[
\llbracket\varphi\rrbracket=\{\,s\colon X\to M\mid M\models\varphi[s]\,\},
\]
és a dir, l'\emph{extensió} d'alguna fórmula. Totes dues són functorials en el context: un canvi de variables $f\colon X\to Y$ transforma les fórmules per substitució ---reemplaçant cada variable de $X$ per la corresponent de $Y$--- i transforma els predicats transportant les assignacions ($s\colon Y\to M$ es restringeix a $s\comp f\colon X\to M$). A cada context $X$ li associem l'aplicació que interpreta cada fórmula com la seva extensió,
\[
\eta_X\colon F(X)\to G(X),
\qquad
\eta_X(\varphi)=\llbracket\varphi\rrbracket.
\]
La condició de naturalitat, per a un canvi de variables $f\colon X\to Y$, és que interpretar la fórmula substituïda coincideixi amb transportar l'extensió de la fórmula original:
\[
\llbracket\varphi[f]\rrbracket=\{\,s\colon Y\to M\mid s\comp f\in\llbracket\varphi\rrbracket\,\}.
\]
Això és exactament el \textbf{lema de substitució} de la teoria de models: el significat d'una fórmula després de canviar-ne les variables es pot calcular transportant el significat original. Per tant $\eta\colon F\Rightarrow G$ és una transformació natural, i expressa que la interpretació és compatible amb la substitució. Com en el cas proposicional de la secció~\ref{sec:cap-transformacio}, $\eta_X$ no és injectiva ---dues fórmules lògicament equivalents defineixen el mateix predicat---, de manera que la naturalitat no és aquí una identificació, sinó el principi que \emph{el significat d'una fórmula no depèn de com n'hàgim anomenat les variables}.

Aquest exemple demana una convenció essencial: la substitució ha de ser \emph{lliure de captura}, i cal fixar el tractament de les variables lligades. Si les fórmules es consideren cadenes literals i es trien variables noves sense cap identificació posterior, les lleis estrictes del functor no queden justificades. Per exemple, en $\varphi=\exists y\, R(x,y)$, substituir la variable lliure $x$ per $y$ no pot donar ingènuament $\exists y\, R(y,y)$: cal reanomenar primer la variable lligada, i obtenir, per exemple, $\exists z\, R(y,z)$. Adoptem, doncs, la convenció següent: disposem d'una reserva infinita de variables; treballem amb fórmules mòdul reanomenament de variables lligades, i totes les substitucions s'entenen lliures de captura; els contextos són conjunts finits de variables. Amb aquesta convenció, les substitucions induïdes per aplicacions entre contextos respecten identitats i composició.

Convé també indicar explícitament que $F,G$ prenen valors a $\cat{Set}$, i escriure l'acció de $G$ sobre un canvi de variables $f\colon X\to Y$ com
\[
G(f)(S)=\{\,s\in M^Y\mid s\comp f\in S\,\},\qquad S\subseteq M^X.
\]
Això mostra, a més, que la imatge inversa d'un predicat definible és definible, pel lema de substitució. Finalment, en una estructura $M$ fixada dues fórmules poden tenir la mateixa extensió sense ser lògicament equivalents en totes les estructures; és útil distingir aquesta \emph{equivalència semàntica relativa a $M$} de la \emph{validesa lògica} general.
\end{exemple}

\begin{exemple}[Doble dual d'espais vectorials]\index{transformació natural!doble dual}
Fixat un cos $\mathbb K$, hi ha el functor identitat sobre els espais vectorials i el functor \textbf{doble dual} $V\mapsto V^{**}$. L'aplicació que envia cada vector $v$ a l'avaluació \enquote{avaluar en $v$},
\[
\eta_V\colon V\to V^{**},
\qquad
\eta_V(v)(\varphi)=\varphi(v),
\]
és una transformació natural de la identitat cap al doble dual. En aquest cas, la naturalitat es pot llegir dient que la identificació de $V$ amb la seva imatge dins $V^{**}$ no fa servir cap elecció de base: la construcció $v\mapsto(\varphi\mapsto\varphi(v))$ és intrínseca. Aquesta lectura ---\enquote{independent de la base}--- és una interpretació d'aquest exemple concret, no la definició de naturalitat, que continua sent la commutativitat dels quadrats. En restringir-se als espais de dimensió finita, cada $\eta_V$ és un isomorfisme, cosa que ens porta a la noció següent.
\end{exemple}

\section{Isomorfismes naturals}

\begin{definicio}\index{isomorfisme natural}
Una transformació natural $\eta\colon F\Rightarrow G$ és un \textbf{isomorfisme natural} si cada component $\eta_A$ és un isomorfisme a $\cat{D}$. En aquest cas escrivim $F\cong G$ i diem que $F$ i $G$ són \textbf{naturalment isomorfs}.
\end{definicio}

\begin{proposicio}
Si $\eta\colon F\Rightarrow G$ és un isomorfisme natural, aleshores els inversos $\eta_A^{-1}$ són els components d'una transformació natural $\eta^{-1}\colon G\Rightarrow F$.
\end{proposicio}

\begin{prova}
Cal comprovar que els $\eta_A^{-1}\colon G(A)\to F(A)$ satisfan la condició de naturalitat. Partim del quadrat de $\eta$ per a un morfisme $f\colon A\to B$:
\[
G(f)\comp\eta_A=\eta_B\comp F(f).
\]
Composant per l'esquerra amb $\eta_B^{-1}$ i per la dreta amb $\eta_A^{-1}$,
\[
\eta_B^{-1}\comp G(f)=\eta_B^{-1}\comp G(f)\comp\eta_A\comp\eta_A^{-1}=\eta_B^{-1}\comp\eta_B\comp F(f)\comp\eta_A^{-1}=F(f)\comp\eta_A^{-1}.
\]
Aquesta és la condició de naturalitat per a $\eta^{-1}$. Per tant $\eta^{-1}\colon G\Rightarrow F$ és una transformació natural, i és inversa de $\eta$ component a component.
\end{prova}

\begin{observacio}\label{obs:iso-no-natural}
La distinció entre isomorfisme i isomorfisme \emph{natural} és fonamental. Que $F(A)$ i $G(A)$ siguin isomorfs per a cada $A$ no basta: cal que els isomorfismes es puguin triar \emph{alhora} de manera compatible amb tots els morfismes.

Un exemple mínim ho mostra sense cap tecnicisme. Prenem el grup cíclic d'ordre dos, $\mathbb{Z}/2\mathbb{Z}=\{1,g\}$ amb $g^2=1$, vist com a categoria d'un sol objecte $\cat{B}(\mathbb{Z}/2\mathbb{Z})$, i dos functors $F,G\colon\cat{B}(\mathbb{Z}/2\mathbb{Z})\to\cat{Set}$. Tots dos envien l'únic objecte $\ast$ al mateix conjunt $\{0,1\}$, però el generador $g$ actua de manera diferent: en $F$ actua per la identitat, i en $G$ per la transposició $\tau$ que intercanvia $0$ i $1$. Puntualment no hi ha res a distingir: $F(\ast)=G(\ast)=\{0,1\}$ són el mateix conjunt, i per tant \emph{isomorfs}. Tanmateix, no hi ha cap isomorfisme natural $F\cong G$. Una transformació natural $\eta\colon F\Rightarrow G$ tindria una única component $\eta_\ast\colon\{0,1\}\to\{0,1\}$, i la naturalitat respecte del morfisme $g$ exigiria
\[
G(g)\comp\eta_\ast=\eta_\ast\comp F(g),
\qquad\text{és a dir}\qquad
\tau\comp\eta_\ast=\eta_\ast.
\]
Però $\tau$ no té punts fixos, de manera que cap aplicació $\eta_\ast$ pot complir $\tau\comp\eta_\ast=\eta_\ast$ (aplicant a qualsevol $x$ obtindríem $\tau(\eta_\ast(x))=\eta_\ast(x)$, impossible). No hi ha, doncs, cap transformació natural $F\Rightarrow G$, i menys encara un isomorfisme natural: la igualtat puntual dels objectes no es pot fer compatible amb l'acció del morfisme.

El cas clàssic del dual d'un espai vectorial il·lustra la mateixa idea, però demana una precaució addicional. Sovint es diu que un espai de dimensió finita \enquote{és isomorf al seu dual $V\cong V^*$, però no naturalment}. Cal formular-ho amb cura, perquè la identitat és un functor \emph{covariant} $\cat{Vect}_{\mathbb K}\to\cat{Vect}_{\mathbb K}$ mentre que el dual $(-)^*$ és \emph{contravariant}, $\cat{Vect}_{\mathbb K}\op\to\cat{Vect}_{\mathbb K}$: no són dos functors paral·lels als quals es pugui aplicar directament la definició de transformació natural d'aquest capítol. La manera correcta de precisar-ho és restringir-se als \emph{isomorfismes} lineals i comparar la identitat amb l'acció covariant $T\mapsto(T^{-1})^*$; aleshores es comprova que cap família d'isomorfismes $V\cong V^*$ és compatible amb tots els isomorfismes lineals. Aquest desenvolupament es fa amb detall a l'exercici~\ref{ex:dual-no-natural} de la part de pràctica. En canvi, l'isomorfisme $V\cong V^{**}$ amb el doble dual sí que està ben plantejat sobre \emph{totes} les aplicacions lineals ---identitat i doble dual són tots dos covariants--- i és natural.
\end{observacio}

\section{La categoria de functors}

Les transformacions naturals es poden compondre, i això organitza els functors entre dues categories en una nova categoria.

\begin{definicio}[Composició vertical]\index{transformació natural!composició vertical}
Donades dues transformacions naturals $\eta\colon F\Rightarrow G$ i $\theta\colon G\Rightarrow H$ entre functors $\cat{C}\to\cat{D}$, la seva \textbf{composició vertical} $\theta\comp\eta\colon F\Rightarrow H$ té per components
\[
(\theta\comp\eta)_A=\theta_A\comp\eta_A.
\]
El nom \emph{vertical} es llegeix apilant els dos quadrats de naturalitat per a un morfisme $f\colon A\to B$, un damunt de l'altre, i component les columnes:
\[
\begin{tikzpicture}[baseline=(m.center)]
  \matrix (m) [matrix of math nodes, row sep=2.6em, column sep=3.4em] {
    F(A) & G(A) & H(A) \\
    F(B) & G(B) & H(B) \\
  };
  \draw[-{Stealth[scale=1.1]}] (m-1-1) -- node[above] {$\eta_A$} (m-1-2);
  \draw[-{Stealth[scale=1.1]}] (m-1-2) -- node[above] {$\theta_A$} (m-1-3);
  \draw[-{Stealth[scale=1.1]}] (m-2-1) -- node[below] {$\eta_B$} (m-2-2);
  \draw[-{Stealth[scale=1.1]}] (m-2-2) -- node[below] {$\theta_B$} (m-2-3);
  \draw[-{Stealth[scale=1.1]}] (m-1-1) -- node[left] {$F(f)$} (m-2-1);
  \draw[-{Stealth[scale=1.1]}] (m-1-2) -- node[left] {$G(f)$} (m-2-2);
  \draw[-{Stealth[scale=1.1]}] (m-1-3) -- node[right] {$H(f)$} (m-2-3);
\end{tikzpicture}
\]
Els dos quadrats commuten (naturalitat de $\eta$ i de $\theta$), per tant el rectangle exterior també, que és la naturalitat de $\theta\comp\eta$.
\end{definicio}

\begin{proposicio}
La composició vertical torna a ser una transformació natural, és associativa, i la transformació identitat $\id_F$ n'és l'element neutre.
\end{proposicio}

\begin{prova}
Per a la naturalitat, considerem un morfisme $f\colon A\to B$. Usant successivament els quadrats de naturalitat de $\eta$ i de $\theta$,
\[
H(f)\comp\theta_A\comp\eta_A=\theta_B\comp G(f)\comp\eta_A=\theta_B\comp\eta_B\comp F(f).
\]
Per tant $\theta\comp\eta$ satisfà la condició de naturalitat. L'associativitat i la propietat de la identitat es redueixen a les propietats corresponents de la composició a $\cat{D}$, aplicades component a component.
\end{prova}

\begin{definicio}\index{categoria!de functors}\index{DC@$\cat{D}^{\cat{C}}$}
Sigui $\cat{C}$ una categoria petita i $\cat{D}$ una categoria qualsevol. La \textbf{categoria de functors} $\cat{D}^{\cat{C}}$ (també escrita $[\cat{C},\cat{D}]$ o $\operatorname{Fun}(\cat{C},\cat{D})$) té:
\begin{itemize}
\item per objectes, els functors $\cat{C}\to\cat{D}$;
\item per morfismes, les transformacions naturals entre ells;
\item per composició, la composició vertical, i per identitats, les transformacions identitat.
\end{itemize}
\end{definicio}

\begin{observacio}
Els isomorfismes de la categoria de functors $\cat{D}^{\cat{C}}$ són exactament els isomorfismes naturals. Així, dir que dos functors són naturalment isomorfs és dir que són isomorfs com a objectes de $\cat{D}^{\cat{C}}$; recuperem la noció d'isomorfisme del primer capítol (la definició~\ref{def:isomorfisme}), ara un nivell més amunt.
\end{observacio}

\begin{observacio}[Qüestions de mida]\index{transformació natural!Nat@$\operatorname{Nat}$}
Convé dir per què la definició demana $\cat{C}$ petita. La raó de fons és la mateixa que impedia formar una categoria de \emph{totes} les categories al capítol anterior. En el marc de classes i conjunts que adoptem (vegeu l'apèndix de convenció de mida), si $\cat{C}$ fos gran, un functor $\cat{C}\to\cat{D}$ podria ser una dada de classe pròpia ---la seva assignació sobre una classe pròpia d'objectes---, i les classes pròpies no poden ser elements de cap col·lecció; aleshores els functors $\cat{C}\to\cat{D}$ no es poden reunir automàticament com els \emph{objectes} d'una categoria ordinària del mateix nivell. Amb $\cat{C}$ petita, en canvi, cada functor és una dada de conjunt i la construcció no té cap problema. Això no impedeix parlar \emph{individualment} de functors i transformacions naturals amb domini gran ---ho fem contínuament, com amb $\cat{Set}\to\cat{Set}$---; el que no fem és agrupar-los tots en una categoria de functors.

Amb $\cat{C}$ petita, $\cat{D}^{\cat{C}}$ pot ser encara gran, però és \textbf{localment petita} tan bon punt $\cat{D}$ ho és: una transformació natural $F\Rightarrow G$ queda determinada pels seus components $\eta_A$, un per a cada objecte $A$ de $\cat{C}$, de manera que la col·lecció de totes les transformacions naturals $F\Rightarrow G$ és un subconjunt del producte
\[
\operatorname{Nat}(F,G)\;\subseteq\;\prod_{A\in\Ob(\cat{C})}\Hom_{\cat{D}}(F(A),G(A)),
\]
indexat pel conjunt d'objectes de $\cat{C}$, i per tant és un conjunt. Aquest \textbf{homconjunt} (de l'anglès \emph{hom-set}; el prefix \enquote{hom-}, de \emph{homomorphism}, no es tradueix) s'escriu sovint
\[
\operatorname{Nat}(F,G)
\]
en lloc del més feixuc $\Hom_{\cat{D}^{\cat{C}}}(F,G)$.
\end{observacio}

\section{Els grafs com a categoria de functors}\index{graf!com a categoria de functors}\index{categoria!de functors!grafs}

Un dels exemples més il·luminatius de categoria de functors mostra que un objecte aparentment aliè a les categories ---un graf dirigit--- \emph{és}, en realitat, un functor, i que els seus morfismes són transformacions naturals. Aquest exemple reprèn el fil dels grafs que hem seguit des del primer capítol i és, a més, el primer indici que categories de functors del tipus $\cat{Set}^{\cat{C}}$ codifiquen estructures matemàtiques familiars.

Considerem la categoria índex $\cat{P}$ amb \textbf{dos objectes} $E$ i $V$ i \textbf{dues fletxes paral·leles} $s,t\colon E\to V$ (a més de les identitats):
\[
\begin{tikzpicture}[baseline=(E.base)]
  \node (E) at (0,0) {$E$};
  \node (V) at (2.4,0) {$V$};
  \draw[-{Stealth[scale=1.1]}] ([yshift=3pt]E.east) -- node[above] {$s$} ([yshift=3pt]V.west);
  \draw[-{Stealth[scale=1.1]}] ([yshift=-3pt]E.east) -- node[below] {$t$} ([yshift=-3pt]V.west);
\end{tikzpicture}
\]

\begin{proposicio}\index{graf!com a categoria de functors}
La categoria de functors $\cat{Set}^{\cat{P}}$ és (isomorfa a) la categoria $\cat{Graph}$ dels grafs dirigits i els morfismes de grafs.
\end{proposicio}

\begin{prova}
Un functor $G\colon\cat{P}\to\cat{Set}$ consisteix en dos conjunts $G(E)$ i $G(V)$ i dues aplicacions $G(s),G(t)\colon G(E)\to G(V)$. Interpretant $G(E)$ com el conjunt d'\textbf{arestes}, $G(V)$ com el conjunt de \textbf{vèrtexs} i $G(s),G(t)$ com les aplicacions \enquote{origen} i \enquote{destí}, aquesta dada és \emph{exactament} un graf dirigit. Recíprocament, tot graf dirigit determina un tal functor. Els objectes de $\cat{Set}^{\cat{P}}$ són, doncs, els grafs dirigits.

Vegem ara els morfismes. Una transformació natural $\varphi\colon G\Rightarrow H$ entre dos d'aquests functors té dos components, $\varphi_E\colon G(E)\to H(E)$ i $\varphi_V\colon G(V)\to H(V)$, i la condició de naturalitat per a les fletxes $s$ i $t$ dona els dos quadrats
\[
\begin{tikzpicture}[baseline=(m.center)]
  \matrix (m) [matrix of math nodes, row sep=2.4em, column sep=3em] {
    G(E) & H(E) \\
    G(V) & H(V) \\
  };
  \draw[-{Stealth[scale=1.1]}] (m-1-1) -- node[above] {$\varphi_E$} (m-1-2);
  \draw[-{Stealth[scale=1.1]}] (m-1-1) -- node[left] {$G(s)$} (m-2-1);
  \draw[-{Stealth[scale=1.1]}] (m-1-2) -- node[right] {$H(s)$} (m-2-2);
  \draw[-{Stealth[scale=1.1]}] (m-2-1) -- node[below] {$\varphi_V$} (m-2-2);
\end{tikzpicture}
\qquad\text{i el mateix amb } t.
\]
Aquests quadrats diuen que $\varphi_E$ i $\varphi_V$ respecten l'origen i el destí de cada aresta; és a dir, són precisament la definició de \textbf{morfisme de grafs}. Per tant els morfismes de $\cat{Set}^{\cat{P}}$ són els morfismes de grafs, i $\cat{Set}^{\cat{P}}\cong\cat{Graph}$.
\end{prova}

\begin{observacio}
Aquest exemple té una lectura de fons: un cop decidim representar els grafs com a functors $\cat{P}\to\cat{Set}$, la representació \emph{justifica i recupera canònicament} la definició habitual de morfisme de grafs, perquè els morfismes de $\cat{Set}^{\cat{P}}$ són, per força, les transformacions naturals. És un criteri de coherència que sovint es fa servir en teoria de categories: quan una estructura s'expressa com a categoria de functors, la noció de morfisme que se'n dedueix hauria de coincidir amb la que s'havia adoptat independentment. Les categories del tipus $\cat{Set}^{\cat{C}}$, que codifiquen estructures com els grafs, seran fonamentals quan estudiem els prefeixos i, més endavant, la semàntica categòrica.
\end{observacio}

\section{Composició horitzontal}\index{transformació natural!composició horitzontal}

La composició vertical apila transformacions entre functors \emph{amb el mateix domini i codomini}. Hi ha una segona manera de compondre, que combina transformacions naturals amb functors i, més generalment, transformacions naturals \enquote{de costat}.

Comencem pel cas mixt, sovint anomenat \emph{whiskering}\footnote{El nom ve de l'anglès \emph{whisker}, \enquote{bigoti}. Si es representa la transformació natural $\eta$ com una cel·la entre dos functors en un diagrama, compondre-la amb un functor equival a afegir una fletxa que sobresurt per un dels costats, com un bigoti enganxat a $\eta$. La imatge del bigoti és també l'origen dels qualificatius \enquote{per l'esquerra} i \enquote{per la dreta}, segons el costat pel qual s'afegeix el functor; convé advertir que aquesta assignació d'\enquote{esquerra} i \enquote{dreta} no és uniforme a la bibliografia. Aquí la fixem per la posició del functor en la notació: $K\eta$ (functor a l'esquerra) és la composició per l'esquerra, i $\eta H$ (functor a la dreta), la composició per la dreta.}, en què composem una transformació natural amb un functor.

\begin{definicio}[Composició amb un functor]\index{whiskering}
Sigui $\eta\colon F\Rightarrow G$ una transformació natural entre functors $F,G\colon\cat{C}\to\cat{D}$.
\begin{itemize}
\item Si $K\colon\cat{D}\to\cat{E}$ és un functor, la \textbf{composició per l'esquerra} $K\eta\colon K\comp F\Rightarrow K\comp G$ té per components $(K\eta)_A=K(\eta_A)$.
\item Si $H\colon\cat{B}\to\cat{C}$ és un functor, la \textbf{composició per la dreta} $\eta H\colon F\comp H\Rightarrow G\comp H$ té per components $(\eta H)_B=\eta_{H(B)}$.
\end{itemize}
\end{definicio}

\begin{proposicio}
Les dues construccions anteriors són transformacions naturals.
\end{proposicio}

\begin{prova}
Per a $K\eta$, apliquem el functor $K$ al quadrat de naturalitat de $\eta$ per a un morfisme $f\colon A\to B$: de $G(f)\comp\eta_A=\eta_B\comp F(f)$, i com que $K$ preserva la composició,
\[
K(G(f))\comp K(\eta_A)=K(\eta_B)\comp K(F(f)),
\]
que és la condició de naturalitat per a $K\eta$ respecte de $K\comp F$ i $K\comp G$. Per a $\eta H$, la naturalitat respecte d'un morfisme $b\colon B\to B'$ de $\cat{B}$ és el quadrat de naturalitat de $\eta$ per al morfisme $H(b)$ de $\cat{C}$.
\end{prova}

Amb aquestes dues operacions podem definir la composició horitzontal general.

\begin{definicio}[Composició horitzontal]\index{transformació natural!composició horitzontal}
Siguin $\eta\colon F\Rightarrow G$ (amb $F,G\colon\cat{C}\to\cat{D}$) i $\theta\colon K\Rightarrow L$ (amb $K,L\colon\cat{D}\to\cat{E}$). La seva \textbf{composició horitzontal} $\theta\ast\eta\colon K\comp F\Rightarrow L\comp G$ té per components
\[
(\theta\ast\eta)_A=\theta_{G(A)}\comp K(\eta_A)=L(\eta_A)\comp\theta_{F(A)}.
\]
\end{definicio}

\begin{observacio}
Les dues expressions de $(\theta\ast\eta)_A$ coincideixen precisament pel quadrat de naturalitat de $\theta$ aplicat al morfisme $\eta_A\colon F(A)\to G(A)$. La composició horitzontal es pot llegir com \enquote{primer un whiskering i després l'altre}, en qualsevol dels dos ordres:
\[
\theta\ast\eta=(\theta G)\comp(K\eta)=(L\eta)\comp(\theta F).
\]
La composició vertical i l'horitzontal, juntes, satisfan una llei d'intercanvi que fa de $\cat{Cat}$ una \emph{2-categoria}; no en necessitarem els detalls, però és el marc en què les transformacions naturals adquireixen tota la seva potència.
\end{observacio}

\section{Equivalència de categories}

Quan volem dir que dues categories són \enquote{essencialment la mateixa}, la igualtat estricta és massa rígida i l'isomorfisme de categories, que demana functors inversos exactes ($G\comp F=\id$), també ho és sovint. La noció adequada relaxa aquesta igualtat a un isomorfisme natural: en lloc d'un invers exacte demanarem un \textbf{quasiinvers}, un functor $G$ amb $G\comp F\cong\id$ i $F\comp G\cong\id$.

\begin{definicio}\index{equivalència de categories}
\label{def:equivalencia}
Una \textbf{equivalència de categories} entre $\cat{C}$ i $\cat{D}$ és una parella de functors
\[
F\colon\cat{C}\to\cat{D},
\qquad
G\colon\cat{D}\to\cat{C},
\]
juntament amb isomorfismes naturals
\[
G\comp F\cong\id_{\cat{C}}
\qquad\text{i}\qquad
F\comp G\cong\id_{\cat{D}}.
\]
Quan existeix una tal equivalència, diem que $\cat{C}$ i $\cat{D}$ són \textbf{equivalents} i escrivim $\cat{C}\simeq\cat{D}$.
\end{definicio}

\begin{observacio}
La diferència amb l'isomorfisme de categories és que aquí no demanem $G\comp F=\id_{\cat{C}}$, sinó només $G\comp F\cong\id_{\cat{C}}$: en tornar de $\cat{C}$ a $\cat{C}$ no recuperem exactament cada objecte, sinó un objecte \emph{naturalment isomorf}. Aquesta flexibilitat és la que fa que l'equivalència sigui la noció \enquote{correcta} de comparació entre categories: respecta tota l'estructura categòrica, que només distingeix els objectes llevat d'isomorfisme.
\end{observacio}

El resultat central caracteritza les equivalències mitjançant les tres propietats dels functors que vam introduir al capítol anterior: plenitud i fidelitat (la definició~\ref{def:ple-fidel}) i l'exhaustivitat essencial (la definició~\ref{def:ess-exhaustiu}).

\begin{teorema}\index{equivalència de categories!caracterització}
\label{teo:caract-equivalencia}
Un functor $F\colon\cat{C}\to\cat{D}$ forma part d'una equivalència de categories si i només si és \textbf{ple}, \textbf{fidel} i \textbf{essencialment exhaustiu}.
\end{teorema}

\begin{prova}[Esquema, amb el nucli de la construcció]
Suposem primer que $F$ forma part d'una equivalència, amb functor $G\colon\cat{D}\to\cat{C}$ i isomorfismes naturals $\eta\colon\id_{\cat{C}}\cong G\comp F$ i $\varepsilon\colon F\comp G\cong\id_{\cat{D}}$. Aleshores $F$ és essencialment exhaustiu, perquè cada $D$ compleix $F(G(D))\cong D$ via $\varepsilon_D$. I és ple i fidel: donada una fletxa $u\colon F(A)\to F(B)$, la composició
\[
A\xrightarrow{\eta_A} G(F(A))\xrightarrow{G(u)} G(F(B))\xrightarrow{\eta_B^{-1}} B
\]
és l'única fletxa $A\to B$ que $F$ envia a $u$ (unicitat per la naturalitat de $\eta$ i la injectivitat que se'n dedueix), de manera que cada $F_{A,B}$ és bijectiva.

Recíprocament, sigui $F$ ple, fidel i essencialment exhaustiu; construïm un \textbf{quasiinvers} $G$ ---un functor amb $G\comp F\cong\id_{\cat{C}}$ i $F\comp G\cong\id_{\cat{D}}$, no un invers exacte. Per l'exhaustivitat essencial, per a cada objecte $D$ de $\cat{D}$ triem un objecte $G(D)$ de $\cat{C}$ i un isomorfisme $\varepsilon_D\colon F(G(D))\xrightarrow{\cong} D$. Per definir $G$ sobre un morfisme $v\colon D\to D'$, considerem la fletxa $\varepsilon_{D'}^{-1}\comp v\comp\varepsilon_D\colon F(G(D))\to F(G(D'))$; com que $F$ és ple i fidel, hi ha una \emph{única} fletxa $G(v)\colon G(D)\to G(D')$ amb
\[
F(G(v))=\varepsilon_{D'}^{-1}\comp v\comp\varepsilon_D.
\]
La unicitat (fidelitat) força $G(\id_D)=\id_{G(D)}$ i $G(v'\comp v)=G(v')\comp G(v)$, de manera que $G$ és un functor. Per construcció, els $\varepsilon_D$ són els components d'un isomorfisme natural $F\comp G\cong\id_{\cat{D}}$; i la plenitud i la fidelitat de $F$ permeten construir, dualment, l'isomorfisme natural $\id_{\cat{C}}\cong G\comp F$. Així $F$ i $G$ formen una equivalència.
\end{prova}

\begin{observacio}[Sobre l'elecció i la remissió als detalls]
La tria simultània d'un objecte $G(D)$ i d'un isomorfisme $\varepsilon_D$ per a cada $D$ fa servir una forma de l'\emph{axioma de l'elecció}. Convé precisar de quina forma: la família triada està indexada pels objectes de $\cat{D}$, i si $\cat{D}$ és gran aquest índex és una classe pròpia, de manera que no n'hi ha prou amb l'elecció per a famílies indexades per un conjunt, sinó que cal una forma d'\emph{elecció global} (disponible, per exemple, a NBG amb l'axioma d'elecció global; vegeu l'apèndix de convenció de mida). En casos concrets sovint se'n pot triar un quasiinvers canònic sense recórrer-hi. La verificació completa de les naturalitats ---els dos quadrats que fan de $\eta$ i $\varepsilon$ isomorfismes naturals--- es desplega pas a pas a l'exercici resolt~\ref{ex:caract-equivalencia} de la part de Pràctica (secció de la caracterització d'equivalències).
\end{observacio}

\begin{definicio}[Esquelet]\index{categoria!esquelètica}\index{esquelet}\label{def:esquelet}
Una categoria és \textbf{esquelètica} si dos objectes isomorfs qualssevol són \emph{iguals}; és a dir, si cada classe d'isomorfisme té un únic objecte. Un \textbf{esquelet} d'una categoria $\cat{C}$ és una subcategoria plena que conté \emph{exactament} un objecte de cada classe d'isomorfisme de $\cat{C}$; un esquelet és sempre una categoria esquelètica.
\end{definicio}

\begin{observacio}
Convé distingir dues nocions properes. Una subcategoria plena que conté \emph{almenys} un objecte de cada classe d'isomorfisme ja és equivalent a $\cat{C}$ (com veurem tot seguit), però pot tenir diversos representants isomorfs d'una mateixa classe. Un \emph{esquelet} n'és el cas ajustat: en conté \emph{exactament} un, de manera que no queden objectes isomorfs distints. Tota categoria (localment petita) en té un, i tots els seus esquelets són equivalents entre si.
\end{observacio}

\begin{exemple}
Tota categoria $\cat{C}$ és equivalent a qualsevol subcategoria plena que contingui almenys un objecte de cada classe d'isomorfisme de $\cat{C}$. En efecte, la inclusió d'aquesta subcategoria és plena i fidel per ser una subcategoria plena, i essencialment exhaustiva per contenir un representant de cada classe. Per exemple, la categoria $\cat{FinSet}$ de tots els conjunts finits és equivalent a la seva subcategoria plena formada pels conjunts $\underline{n}=\{1,2,\dots,n\}$ amb $n\geq 0$: no cal treballar amb tots els conjunts finits, sinó només amb un de cada mida. Aquesta subcategoria $\{\underline{n}\mid n\geq 0\}$ és, de fet, un \emph{esquelet} de $\cat{FinSet}$, perquè conté exactament un conjunt de cada cardinalitat finita. Convé notar que $\cat{FinSet}$ és una categoria \emph{gran} ---els seus objectes són massa nombrosos per formar un conjunt, tal com passa amb $\cat{Set}$--- mentre que l'esquelet $\{\underline{n}\mid n\geq 0\}$ és \emph{petit}, ja que té un conjunt (de fet, numerable) d'objectes. L'equivalència mostra, doncs, que $\cat{FinSet}$ és \emph{essencialment petita}.
\end{exemple}

\begin{observacio}
Aquest exemple mostra la utilitat pràctica de l'equivalència: permet substituir una categoria gran i redundant per una de més petita i manejable, sense perdre cap informació categòrica. És el reflex, a nivell de categories, del fet que en teoria de categories els objectes només importen llevat d'isomorfisme.
\end{observacio}

\section{Cap a les propietats universals}

Amb les transformacions naturals hem completat una torre de tres nivells: objectes i morfismes dins d'una categoria, functors entre categories, i transformacions naturals entre functors. Aquesta estructura ---sovint resumida dient que les categories, els functors i les transformacions naturals formen una \emph{2-categoria}--- és el marc en què es formulen les construccions més potents de la teoria.

La primera d'aquestes construccions són les \textbf{propietats universals}, que caracteritzen un objecte no per la seva constitució interna, sinó per la representabilitat d'un functor associat ---normalment amb un element universal--- que el determina de manera única llevat d'isomorfisme. El producte i el coproducte que hem trobat de passada, els objectes inicial i terminal, i els límits i colímits en general, en són casos particulars. Els estudiarem als capítols de representabilitat i de límits, on veurem que aquest llenguatge permet expressar-los amb precisió i deduir-ne les propietats de manera uniforme.

\begin{resumcap}
\apartat{Cinc idees} (1)~Una transformació natural és una família de morfismes que commuta amb l'acció dels dos functors. (2)~Un isomorfisme natural és una transformació natural invertible component a component. (3)~Els functors $\cat{C}\to\cat{D}$ i les seves transformacions formen la categoria $\cat{D}^{\cat{C}}$. (4)~Els grafs són els objectes d'una categoria de functors $\cat{Set}^{\cat{P}}$. (5)~Una equivalència és un functor ple, fidel i essencialment exhaustiu; preserva les nocions definides llevat d'isomorfisme, no les que depenen de la identitat dels objectes.
\apartat{Errors típics} Oblidar comprovar la naturalitat (que el quadrat commuti per a \emph{tot} morfisme); confondre isomorfisme natural amb igualtat de functors; creure que una equivalència conserva propietats com \enquote{tenir un sol objecte}; i barrejar composició vertical i horitzontal de transformacions.
\apartat{Lectura recomanada} \cite[cap.~7--8]{awodey} per a transformacions naturals i equivalències; \cite[cap.~1]{leinster}; \cite[I.4 i IV]{maclane} per a la composició horitzontal.
\end{resumcap}
